\documentclass[11pt]{article}
\usepackage[utf8]{inputenc}
\usepackage{amsmath, amssymb, amsthm}
\usepackage{geometry}
\usepackage{graphicx}
\usepackage{enumitem}
\usepackage{fancyhdr}
\usepackage{xcolor}

\geometry{a4paper, margin=1in}

% Header/Footer
\pagestyle{fancy}
\fancyhf{}
\lhead{Homework 12.4: The Cross Product}
\rhead{Tashfeen Omran}
\cfoot{\thepage}

% Custom commands for vectors to ensure bold notation
\renewcommand{\vec}[1]{\mathbf{#1}}
\newcommand{\ihat}{\mathbf{i}}
\newcommand{\jhat}{\mathbf{j}}
\newcommand{\khat}{\mathbf{k}}

\title{\textbf{Homework 12.4: The Cross Product}}
\author{Tashfeen Omran}
\date{February 2026}

\begin{document}

\maketitle

\tableofcontents
\newpage

\section{Part 1: Comprehensive Introduction, Context, and Prerequisites}

\subsection{Core Concepts}
The \textbf{Cross Product} (or vector product) is a binary operation on two vectors in three-dimensional space, denoted as $\vec{a} \times \vec{b}$. Unlike the dot product, which results in a scalar, the cross product results in a \textbf{vector}.

\textbf{Key Geometric Properties:}
\begin{enumerate}
    \item \textbf{Direction:} The vector $\vec{c} = \vec{a} \times \vec{b}$ is orthogonal (perpendicular) to both $\vec{a}$ and $\vec{b}$. Its direction is determined by the \textbf{Right-Hand Rule}.
    \item \textbf{Magnitude:} The magnitude $|\vec{a} \times \vec{b}|$ represents the area of the parallelogram spanned by $\vec{a}$ and $\vec{b}$.
\end{enumerate}

\subsection{Intuition and Derivation}
While the Dot Product measures "parallelness" (how much two vectors point in the same direction), the Cross Product measures "perpendicularness." 

Derivations often stem from trying to find a vector $\vec{v} = \langle x, y, z \rangle$ such that $\vec{a} \cdot \vec{v} = 0$ and $\vec{b} \cdot \vec{v} = 0$. Solving this system of linear equations leads to the standard determinant formula. The magnitude formula is derived using the identity $|\vec{a} \times \vec{b}|^2 = |\vec{a}|^2|\vec{b}|^2 - (\vec{a} \cdot \vec{b})^2$ (Lagrange's Identity), which simplifies to $|\vec{a}||\vec{b}|\sin \theta$.

\subsection{Historical Context and Motivation}
In the mid-19th century, Sir William Rowan Hamilton developed \textit{Quaternions} to describe mechanics in 3D space. The quaternion product contained a scalar part and a vector part. Later, physicists Josiah Willard Gibbs and Oliver Heaviside separated these parts into the modern Dot Product and Cross Product. The Cross Product was specifically motivated by problems in electromagnetism (calculating magnetic forces, which act perpendicular to velocity and field) and rotational mechanics (torque).

\subsection{Key Formulas}
\textbf{1. Component Definition (Determinant Form):}
\[
\vec{a} \times \vec{b} = \begin{vmatrix} \ihat & \jhat & \khat \\ a_1 & a_2 & a_3 \\ b_1 & b_2 & b_3 \end{vmatrix} = (a_2 b_3 - a_3 b_2)\ihat - (a_1 b_3 - a_3 b_1)\jhat + (a_1 b_2 - a_2 b_1)\khat
\]

\textbf{2. Geometric Definition:}
\[
|\vec{a} \times \vec{b}| = |\vec{a}| |\vec{b}| \sin \theta
\]
where $\theta$ is the angle between the vectors ($0 \le \theta \le \pi$).

\textbf{3. Scalar Triple Product (Volume):}
\[
V = |\vec{a} \cdot (\vec{b} \times \vec{c})|
\]
This calculates the volume of the parallelepiped determined by $\vec{a}, \vec{b}, \text{ and } \vec{c}$.

\subsection{Prerequisites}
\begin{itemize}
    \item \textbf{Determinants:} Ability to calculate $2 \times 2$ and $3 \times 3$ determinants.
    \item \textbf{Vector Algebra:} Vector addition, subtraction, and standard basis vectors $\ihat, \jhat, \khat$.
    \item \textbf{Dot Product:} Understanding orthogonality ($\vec{u} \cdot \vec{v} = 0$).
    \item \textbf{Trigonometry:} Values of sine and cosine for standard angles; Pythagorean identities.
\end{itemize}

\newpage

\section{Part 2: Detailed Homework Solutions}

\subsection*{Problem 1: SCalcET9 12.4.001}
\textbf{Given:} $\vec{a} = \langle 2, 3, 0 \rangle$, $\vec{b} = \langle 1, 0, 5 \rangle$.

\textbf{Step 1: Calculate the Cross Product.}
\[
\vec{a} \times \vec{b} = \begin{vmatrix} \ihat & \jhat & \khat \\ 2 & 3 & 0 \\ 1 & 0 & 5 \end{vmatrix}
\]
Expand along the first row:
\[
= \ihat \begin{vmatrix} 3 & 0 \\ 0 & 5 \end{vmatrix} - \jhat \begin{vmatrix} 2 & 0 \\ 1 & 5 \end{vmatrix} + \khat \begin{vmatrix} 2 & 3 \\ 1 & 0 \end{vmatrix}
\]
Calculate $2 \times 2$ determinants:
\[
= \ihat (3(5) - 0(0)) - \jhat (2(5) - 0(1)) + \khat (2(0) - 3(1))
\]
\[
= \ihat (15) - \jhat (10) + \khat (-3) = \langle 15, -10, -3 \rangle
\]

\textbf{Step 2: Verify Orthogonality.}
\[
(\vec{a} \times \vec{b}) \cdot \vec{a} = \langle 15, -10, -3 \rangle \cdot \langle 2, 3, 0 \rangle = (15)(2) + (-10)(3) + (-3)(0) = 30 - 30 + 0 = 0
\]
\[
(\vec{a} \times \vec{b}) \cdot \vec{b} = \langle 15, -10, -3 \rangle \cdot \langle 1, 0, 5 \rangle = (15)(1) + (-10)(0) + (-3)(5) = 15 + 0 - 15 = 0
\]

\textbf{Final Answer:}
\begin{align*}
\vec{a} \times \vec{b} &= \langle 15, -10, -3 \rangle \\
(\vec{a} \times \vec{b}) \cdot \vec{a} &= 0 \\
(\vec{a} \times \vec{b}) \cdot \vec{b} &= 0
\end{align*}

\hrulefill

\subsection*{Problem 2: SCalcET9 12.4.002}
\textbf{Given:} $\vec{a} = \langle 5, 4, -3 \rangle$, $\vec{b} = \langle 2, -1, 1 \rangle$.

\textbf{Step 1: Calculate Cross Product.}
\[
\vec{a} \times \vec{b} = \begin{vmatrix} \ihat & \jhat & \khat \\ 5 & 4 & -3 \\ 2 & -1 & 1 \end{vmatrix}
\]
\[
= \ihat (4(1) - (-3)(-1)) - \jhat (5(1) - (-3)(2)) + \khat (5(-1) - 4(2))
\]
\[
= \ihat (4 - 3) - \jhat (5 - (-6)) + \khat (-5 - 8)
\]
\[
= 1\ihat - 11\jhat - 13\khat = \langle 1, -11, -13 \rangle
\]

\textbf{Step 2: Verify Orthogonality.}
\[
(\vec{a} \times \vec{b}) \cdot \vec{a} = 1(5) + (-11)(4) + (-13)(-3) = 5 - 44 + 39 = 0
\]
\[
(\vec{a} \times \vec{b}) \cdot \vec{b} = 1(2) + (-11)(-1) + (-13)(1) = 2 + 11 - 13 = 0
\]

\textbf{Final Answer:}
\begin{align*}
\vec{a} \times \vec{b} &= \langle 1, -11, -13 \rangle \\
(\vec{a} \times \vec{b}) \cdot \vec{a} &= 0 \\
(\vec{a} \times \vec{b}) \cdot \vec{b} &= 0
\end{align*}

\hrulefill

\subsection*{Problem 3: SCalcET9 12.4.003}
\textbf{Given:} $\vec{a} = 4\jhat - 9\khat = \langle 0, 4, -9 \rangle$, $\vec{b} = -\ihat + 3\jhat + \khat = \langle -1, 3, 1 \rangle$.

\textbf{Step 1: Calculate Cross Product.}
\[
\vec{a} \times \vec{b} = \begin{vmatrix} \ihat & \jhat & \khat \\ 0 & 4 & -9 \\ -1 & 3 & 1 \end{vmatrix}
\]
\[
= \ihat (4(1) - (-9)(3)) - \jhat (0(1) - (-9)(-1)) + \khat (0(3) - 4(-1))
\]
\[
= \ihat (4 + 27) - \jhat (0 - 9) + \khat (0 + 4)
\]
\[
= 31\ihat + 9\jhat + 4\khat = \langle 31, 9, 4 \rangle
\]

\textbf{Step 2: Verify Orthogonality.}
\[
(\vec{a} \times \vec{b}) \cdot \vec{a} = 31(0) + 9(4) + 4(-9) = 0 + 36 - 36 = 0
\]
\[
(\vec{a} \times \vec{b}) \cdot \vec{b} = 31(-1) + 9(3) + 4(1) = -31 + 27 + 4 = 0
\]

\textbf{Final Answer:}
\begin{align*}
\vec{a} \times \vec{b} &= 31\ihat + 9\jhat + 4\khat \\
(\vec{a} \times \vec{b}) \cdot \vec{a} &= 0 \\
(\vec{a} \times \vec{b}) \cdot \vec{b} &= 0
\end{align*}

\hrulefill

\subsection*{Problem 4: SCalcET9 12.4.006}
\textbf{Given:} $\vec{a} = t\ihat + \cos(t)\jhat + \sin(t)\khat$, $\vec{b} = \ihat - \sin(t)\jhat + \cos(t)\khat$.

\textbf{Step 1: Calculate Cross Product.}
\[
\vec{a} \times \vec{b} = \begin{vmatrix} \ihat & \jhat & \khat \\ t & \cos t & \sin t \\ 1 & -\sin t & \cos t \end{vmatrix}
\]
\textbf{i-component:}
\[
(\cos t)(\cos t) - (\sin t)(-\sin t) = \cos^2 t + \sin^2 t = 1
\]
\textbf{j-component (include negative sign):}
\[
-[ (t)(\cos t) - (\sin t)(1) ] = -[ t\cos t - \sin t ] = \sin t - t\cos t
\]
\textbf{k-component:}
\[
(t)(-\sin t) - (\cos t)(1) = -t\sin t - \cos t
\]

\textbf{Step 2: Verify Orthogonality.}
Dot with $\vec{a}$:
\[
1(t) + (\sin t - t\cos t)(\cos t) + (-t\sin t - \cos t)(\sin t)
\]
\[
= t + \sin t \cos t - t \cos^2 t - t \sin^2 t - \sin t \cos t
\]
\[
= t - t(\cos^2 t + \sin^2 t) = t - t(1) = 0
\]
Dot with $\vec{b}$:
\[
1(1) + (\sin t - t\cos t)(-\sin t) + (-t\sin t - \cos t)(\cos t)
\]
\[
= 1 - \sin^2 t + t \sin t \cos t - t \sin t \cos t - \cos^2 t
\]
\[
= 1 - (\sin^2 t + \cos^2 t) = 1 - 1 = 0
\]

\textbf{Final Answer:}
\begin{align*}
\vec{a} \times \vec{b} &= \ihat + (\sin t - t\cos t)\jhat + (-t\sin t - \cos t)\khat \\
(\vec{a} \times \vec{b}) \cdot \vec{a} &= 0 \\
(\vec{a} \times \vec{b}) \cdot \vec{b} &= 0
\end{align*}

\hrulefill

\subsection*{Problem 5: SCalcET9 12.4.007}
\textbf{Given:} $\vec{a} = \langle t^3, t^2, t \rangle$, $\vec{b} = \langle t, 3t, 9t \rangle$.

\textbf{Step 1: Calculate Cross Product.}
\[
\vec{a} \times \vec{b} = \begin{vmatrix} \ihat & \jhat & \khat \\ t^3 & t^2 & t \\ t & 3t & 9t \end{vmatrix}
\]
\textbf{i-component:} $t^2(9t) - t(3t) = 9t^3 - 3t^2$
\textbf{j-component:} $-(t^3(9t) - t(t)) = -(9t^4 - t^2) = t^2 - 9t^4$
\textbf{k-component:} $t^3(3t) - t^2(t) = 3t^4 - t^3$

\textbf{Step 2: Verify Orthogonality.}
Dot with $\vec{a}$:
\[
t^3(9t^3 - 3t^2) + t^2(t^2 - 9t^4) + t(3t^4 - t^3)
\]
\[
= 9t^6 - 3t^5 + t^4 - 9t^6 + 3t^5 - t^4 = 0
\]
Dot with $\vec{b}$:
\[
t(9t^3 - 3t^2) + 3t(t^2 - 9t^4) + 9t(3t^4 - t^3)
\]
\[
= 9t^4 - 3t^3 + 3t^3 - 27t^5 + 27t^5 - 9t^4 = 0
\]

\textbf{Final Answer:}
\begin{align*}
\vec{a} \times \vec{b} &= \langle 9t^3 - 3t^2, \ t^2 - 9t^4, \ 3t^4 - t^3 \rangle \\
(\vec{a} \times \vec{b}) \cdot \vec{a} &= 0 \\
(\vec{a} \times \vec{b}) \cdot \vec{b} &= 0
\end{align*}

\hrulefill

\subsection*{Problem 6: SCalcET9 12.4.011}
\textbf{Find:} $(\jhat - \khat) \times (\khat - \ihat)$ using properties.

\textbf{Solution:}
Use distributivity: $(\vec{u} + \vec{v}) \times \vec{w} = \vec{u} \times \vec{w} + \vec{v} \times \vec{w}$.
\[
(\jhat - \khat) \times (\khat - \ihat) = \jhat \times (\khat - \ihat) - \khat \times (\khat - \ihat)
\]
\[
= (\jhat \times \khat) - (\jhat \times \ihat) - (\khat \times \khat) + (\khat \times \ihat)
\]
Recall cycle: $\ihat \to \jhat \to \khat \to \ihat$ is positive.
\begin{itemize}
    \item $\jhat \times \khat = \ihat$
    \item $\jhat \times \ihat = -\khat$
    \item $\khat \times \khat = \vec{0}$
    \item $\khat \times \ihat = \jhat$
\end{itemize}
Substitute:
\[
= \ihat - (-\khat) - \vec{0} + \jhat = \ihat + \khat + \jhat
\]
\textbf{Final Answer:} $\ihat + \jhat + \khat$

\hrulefill

\subsection*{Problem 7: SCalcET9 12.4.012}
\textbf{Find:} $(\ihat + \jhat) \times (\ihat - \jhat)$.

\textbf{Solution:}
\[
= \ihat \times (\ihat - \jhat) + \jhat \times (\ihat - \jhat)
\]
\[
= (\ihat \times \ihat) - (\ihat \times \jhat) + (\jhat \times \ihat) - (\jhat \times \jhat)
\]
Recall $\vec{v} \times \vec{v} = \vec{0}$.
\[
= \vec{0} - (\khat) + (-\khat) - \vec{0} = -2\khat
\]
\textbf{Final Answer:} $-2\khat$

\hrulefill

\subsection*{Problem 8: SCalcET9 12.4.013}
\textbf{Solution:}
\begin{enumerate}[label=(\alph*)]
    \item $\vec{a} \cdot (\vec{b} \times \vec{c})$: $\vec{b} \times \vec{c}$ is a vector. Vector dot vector is a scalar. \textbf{Meaningful; Scalar.}
    \item $\vec{a} \times (\vec{b} \cdot \vec{c})$: $\vec{b} \cdot \vec{c}$ is a scalar. We cannot cross a vector with a scalar. \textbf{Meaningless.}
    \item $\vec{a} \times (\vec{b} \times \vec{c})$: $\vec{b} \times \vec{c}$ is a vector. Vector cross vector is a vector. \textbf{Meaningful; Vector.}
    \item $\vec{a} \cdot (\vec{b} \cdot \vec{c})$: $\vec{b} \cdot \vec{c}$ is a scalar. We cannot dot a vector with a scalar. \textbf{Meaningless.}
    \item $(\vec{a} \cdot \vec{b}) \times (\vec{c} \cdot \vec{d})$: Dot products are scalars. Scalar cross scalar is undefined. \textbf{Meaningless.}
    \item $(\vec{a} \times \vec{b}) \cdot (\vec{c} \times \vec{d})$: Both parentheses yield vectors. Vector dot vector is a scalar. \textbf{Meaningful; Scalar.}
\end{enumerate}

\hrulefill

\subsection*{Problem 9: SCalcET9 12.4.014}
\textbf{Given:} $|\vec{u}| = 4, |\vec{v}| = 7, \theta = 60^\circ$.

\textbf{Magnitude:}
\[
|\vec{u} \times \vec{v}| = |\vec{u}| |\vec{v}| \sin \theta = 4 \cdot 7 \cdot \sin(60^\circ) = 28 \left( \frac{\sqrt{3}}{2} \right) = 14\sqrt{3}
\]

\textbf{Direction:}
Using the Right-Hand Rule: Point fingers along $\vec{u}$ (horizontal right), curl toward $\vec{v}$ (down-right). The thumb points \textbf{into the screen}.

\textbf{Final Answer:}
$|\vec{u} \times \vec{v}| = 14\sqrt{3}$. Direction: \textbf{into the screen}.

\hrulefill

\subsection*{Problem 10: SCalcET9 12.4.016}
\textbf{Given:} $\vec{a}$ in xy-plane (y negative, x positive), $\vec{b}$ in direction of $\khat$. $|\vec{a}|=4, |\vec{b}|=3$.
Since $\vec{a}$ is in the xy-plane and $\vec{b}$ is along the z-axis, they are perpendicular. $\theta = 90^\circ$.

\textbf{(a) Magnitude:}
\[
|\vec{a} \times \vec{b}| = |\vec{a}| |\vec{b}| \sin(90^\circ) = 4 \cdot 3 \cdot 1 = 12
\]

\textbf{(b) Components (RHR):}
$\vec{a}$ points roughly in Direction $(1, -1, 0)$. $\vec{b}$ points in $(0, 0, 1)$.
$\vec{a} \times \vec{b}$: Fingers along $\vec{a}$ (South-East), curl up toward $z$ ($\vec{b}$). Thumb points roughly North-East (positive x, positive y).
Let's verify with determinants roughly. $\vec{a} \approx \langle +, -, 0 \rangle$, $\vec{b} = \langle 0, 0, + \rangle$.
$\ihat: (-)(+) - 0 = \text{negative}$. Wait.
Let's use RHR strictly. $\vec{a}$ is in quadrant IV of xy plane. $\vec{b}$ is straight up.
Rotate $\vec{a}$ 90 degrees toward $\vec{b}$? No, cross product is perpendicular to both.
Perpendicular to $z$-axis means it lies in the xy-plane ($z$-comp is 0).
Perpendicular to $\vec{a}$ (which is Q4). The vector perpendicular to a Q4 vector in the xy plane is either Q1 or Q3.
RHR: Fingers Q4, curl Up. Thumb points to Q1 (Positive x, Positive y).
Wait, let's test specific vectors. $\vec{a} = \langle 1, -1, 0 \rangle, \vec{b} = \langle 0, 0, 1 \rangle$.
$\vec{a} \times \vec{b} = \langle (-1)(1)-0, 0-(1)(1), 0 \rangle = \langle -1, -1, 0 \rangle$.
Let's re-evaluate the image description. The image shows $\vec{a}$ in the xy-plane with positive x and positive y? No, the arrow $\vec{a}$ starts at origin and points towards positive x and positive y?
Looking at the crop: The axis labeled $x$ comes out left. $y$ goes right. $z$ goes up.
Vector $\vec{a}$ is between positive $y$ and positive $x$. It looks like Q1 in xy plane.
Vector $\vec{b}$ is along $z$.
If $\vec{a}$ is Q1 ($+,+$) and $\vec{b}$ is $z$ ($+$).
$\ihat$: $(a_y)(b_z) - 0 = +$.
$\jhat$: $-( (a_x)(b_z) - 0 ) = -$.
$\khat$: $0$.
So x-component positive, y-component negative, z-component 0.
Let's re-read the RHR. Fingers along $\vec{a}$ (Q1), curl to $\vec{b}$ (Up). Thumb points to Q4 (positive x, negative y).

\textbf{Final Answer:}
Magnitude: 12.
x-component: Positive.
y-component: Negative.
z-component: 0.

\hrulefill

\subsection*{Problem 11: SCalcET9 12.4.017}
\textbf{Given:} $\vec{a} = \langle 2, -1, 9 \rangle$, $\vec{b} = \langle 2, 2, 1 \rangle$.

\textbf{Step 1: Calculate $\vec{a} \times \vec{b}$.}
\[
\begin{vmatrix} \ihat & \jhat & \khat \\ 2 & -1 & 9 \\ 2 & 2 & 1 \end{vmatrix} = \ihat(-1 - 18) - \jhat(2 - 18) + \khat(4 - (-2))
\]
\[
= -19\ihat + 16\jhat + 6\khat = \langle -19, 16, 6 \rangle
\]

\textbf{Step 2: Calculate $\vec{b} \times \vec{a}$.}
Property: $\vec{b} \times \vec{a} = -(\vec{a} \times \vec{b})$.
\[
= \langle 19, -16, -6 \rangle
\]

\textbf{Final Answer:} $\vec{a} \times \vec{b} = \langle -19, 16, 6 \rangle$; $\vec{b} \times \vec{a} = \langle 19, -16, -6 \rangle$.

\hrulefill

\subsection*{Problem 12: SCalcET9 12.4.018}
\textbf{Given:} $\vec{a} = \langle 3, 0, 3 \rangle, \vec{b} = \langle 4, 3, -3 \rangle, \vec{c} = \langle 0, 3, 2 \rangle$.

\textbf{Part 1: $\vec{a} \times (\vec{b} \times \vec{c})$}
First $\vec{b} \times \vec{c}$:
\[
\begin{vmatrix} \ihat & \jhat & \khat \\ 4 & 3 & -3 \\ 0 & 3 & 2 \end{vmatrix} = \ihat(6 - (-9)) - \jhat(8 - 0) + \khat(12 - 0) = \langle 15, -8, 12 \rangle
\]
Now $\vec{a} \times \langle 15, -8, 12 \rangle$:
\[
\begin{vmatrix} \ihat & \jhat & \khat \\ 3 & 0 & 3 \\ 15 & -8 & 12 \end{vmatrix} = \ihat(0 - (-24)) - \jhat(36 - 45) + \khat(-24 - 0) = \langle 24, 9, -24 \rangle
\]

\textbf{Part 2: $(\vec{a} \times \vec{b}) \times \vec{c}$}
First $\vec{a} \times \vec{b}$:
\[
\begin{vmatrix} \ihat & \jhat & \khat \\ 3 & 0 & 3 \\ 4 & 3 & -3 \end{vmatrix} = \ihat(0 - 9) - \jhat(-9 - 12) + \khat(9 - 0) = \langle -9, 21, 9 \rangle
\]
Now $\langle -9, 21, 9 \rangle \times \vec{c}$:
\[
\begin{vmatrix} \ihat & \jhat & \khat \\ -9 & 21 & 9 \\ 0 & 3 & 2 \end{vmatrix} = \ihat(42 - 27) - \jhat(-18 - 0) + \khat(-27 - 0) = \langle 15, 18, -27 \rangle
\]

\textbf{Final Answer:} They are not equal.
$\vec{a} \times (\vec{b} \times \vec{c}) = \langle 24, 9, -24 \rangle$.
$(\vec{a} \times \vec{b}) \times \vec{c} = \langle 15, 18, -27 \rangle$.

\hrulefill

\subsection*{Problem 13: SCalcET9 12.4.027}
\textbf{Given:} Vertices $A(-3, 0), B(-1, 5), C(7, 4), D(5, -1)$.
Since these are 2D points, treat $z = 0$.
Vectors defining the parallelogram from $A$:
$\vec{AB} = \langle -1 - (-3), 5 - 0, 0 \rangle = \langle 2, 5, 0 \rangle$.
$\vec{AD} = \langle 5 - (-3), -1 - 0, 0 \rangle = \langle 8, -1, 0 \rangle$.

\textbf{Area:} Magnitude of cross product.
\[
\vec{AB} \times \vec{AD} = \begin{vmatrix} \ihat & \jhat & \khat \\ 2 & 5 & 0 \\ 8 & -1 & 0 \end{vmatrix} = \ihat(0) - \jhat(0) + \khat(2(-1) - 5(8)) = \langle 0, 0, -42 \rangle
\]
Magnitude = $\sqrt{0^2 + 0^2 + (-42)^2} = 42$.
\textbf{Final Answer:} 42.

\hrulefill

\subsection*{Problem 14: SCalcET9 12.4.027.EP}
\textbf{Given:} $A(-3, 0), B(-1, 6), C(8, 5), D(6, -1)$. (Note: coordinates different from Q13).
Treat as 3D points with $z=0$.

\textbf{Step 1: Vectors}
$\vec{AB} = \langle -1 - (-3), 6 - 0, 0 \rangle = \langle 2, 6, 0 \rangle$.
$\vec{AD} = \langle 6 - (-3), -1 - 0, 0 \rangle = \langle 9, -1, 0 \rangle$.

\textbf{Step 2: Cross Product}
\[
\vec{AB} \times \vec{AD} = \begin{vmatrix} \ihat & \jhat & \khat \\ 2 & 6 & 0 \\ 9 & -1 & 0 \end{vmatrix} = \ihat(0) - \jhat(0) + \khat(-2 - 54) = \langle 0, 0, -56 \rangle
\]

\textbf{Step 3: Area}
Area = $|-56| = 56$.

\textbf{Final Answer:}
$\vec{AB} = \langle 2, 6, 0 \rangle$; $\vec{AD} = \langle 9, -1, 0 \rangle$;
$\vec{AB} \times \vec{AD} = \langle 0, 0, -56 \rangle$; Area = 56.

\hrulefill

\subsection*{Problem 15: SCalcET9 12.4.028}
\textbf{Given:} $P(3, 1, 2), Q(6, 2, 5), R(10, 8, 13), S(7, 7, 10)$.
Parallelogram implies $\vec{PQ}$ and $\vec{PS}$ are adjacent sides (check if $\vec{QR} = \vec{PS}$).
$\vec{PQ} = \langle 3, 1, 3 \rangle$.
$\vec{PS} = \langle 4, 6, 8 \rangle$.

\textbf{Cross Product:}
\[
\vec{PQ} \times \vec{PS} = \begin{vmatrix} \ihat & \jhat & \khat \\ 3 & 1 & 3 \\ 4 & 6 & 8 \end{vmatrix}
\]
$\ihat: 8 - 18 = -10$.
$\jhat: -(24 - 12) = -12$.
$\khat: 18 - 4 = 14$.
Vector: $\langle -10, -12, 14 \rangle$.

\textbf{Area:}
$|\langle -10, -12, 14 \rangle| = \sqrt{100 + 144 + 196} = \sqrt{440} = \sqrt{4 \cdot 110} = 2\sqrt{110}$.

\textbf{Final Answer:} $2\sqrt{110}$.

\hrulefill

\subsection*{Problem 16: SCalcET9 12.4.031}
\textbf{Given:} $P(0, -2, 0), Q(5, 1, -2), R(5, 2, 1)$.

\textbf{(a) Orthogonal Vector:} Find $\vec{n} = \vec{PQ} \times \vec{PR}$.
$\vec{PQ} = \langle 5, 3, -2 \rangle$.
$\vec{PR} = \langle 5, 4, 1 \rangle$.
\[
\vec{n} = \begin{vmatrix} \ihat & \jhat & \khat \\ 5 & 3 & -2 \\ 5 & 4 & 1 \end{vmatrix} = \ihat(3 - (-8)) - \jhat(5 - (-10)) + \khat(20 - 15)
\]
\[
= 11\ihat - 15\jhat + 5\khat = \langle 11, -15, 5 \rangle
\]

\textbf{(b) Area of Triangle:}
Area = $\frac{1}{2} |\vec{PQ} \times \vec{PR}| = \frac{1}{2} \sqrt{11^2 + (-15)^2 + 5^2}$.
$= \frac{1}{2} \sqrt{121 + 225 + 25} = \frac{1}{2} \sqrt{371}$.

\textbf{Final Answer:}
Vector: $\langle 11, -15, 5 \rangle$. Area: $\frac{\sqrt{371}}{2}$.

\hrulefill

\subsection*{Problem 17: SCalcET9 12.4.033}
\textbf{Given:} $\vec{a} = \langle 1, 4, 4 \rangle, \vec{b} = \langle -1, 1, 5 \rangle, \vec{c} = \langle 4, 1, 3 \rangle$.
Volume of parallelepiped = $|\vec{a} \cdot (\vec{b} \times \vec{c})|$.

\textbf{Scalar Triple Product:}
\[
\det \begin{pmatrix} 1 & 4 & 4 \\ -1 & 1 & 5 \\ 4 & 1 & 3 \end{pmatrix}
\]
\[
= 1(3 - 5) - 4(-3 - 20) + 4(-1 - 4)
\]
\[
= 1(-2) - 4(-23) + 4(-5) = -2 + 92 - 20 = 70
\]
Volume = $|70| = 70$.

\textbf{Final Answer:} 70 cubic units.

\hrulefill

\subsection*{Problem 18: SCalcET9 12.4.038}
\textbf{Given:} $A(2, 1, 2), B(3, -5, 7), C(6, -2, 0), D(5, 4, -5)$.
Check if coplanar. Compute scalar triple product of vectors from $A$.
$\vec{AB} = \langle 1, -6, 5 \rangle$.
$\vec{AC} = \langle 4, -3, -2 \rangle$.
$\vec{AD} = \langle 3, 3, -7 \rangle$.

\[
\det \begin{pmatrix} 1 & -6 & 5 \\ 4 & -3 & -2 \\ 3 & 3 & -7 \end{pmatrix}
\]
\[
= 1(21 - (-6)) - (-6)(-28 - (-6)) + 5(12 - (-9))
\]
\[
= 1(27) + 6(-22) + 5(21)
\]
\[
= 27 - 132 + 105 = 132 - 132 = 0
\]
Since volume is 0, they are coplanar.

\textbf{Final Answer:} Yes, they lie in the same plane.

\hrulefill

\subsection*{Problem 19: SCalcET9 12.4.043}
\textbf{Given:} $\vec{a} \cdot \vec{b} = 9\sqrt{3}$ and $\vec{a} \times \vec{b} = \langle 3, 6, 6 \rangle$.
$|\vec{a} \times \vec{b}| = \sqrt{3^2 + 6^2 + 6^2} = \sqrt{9 + 36 + 36} = \sqrt{81} = 9$.

We know:
$|\vec{a} \times \vec{b}| = |\vec{a}||\vec{b}| \sin \theta = 9$
$\vec{a} \cdot \vec{b} = |\vec{a}||\vec{b}| \cos \theta = 9\sqrt{3}$

Divide the two equations:
\[
\tan \theta = \frac{9}{9\sqrt{3}} = \frac{1}{\sqrt{3}}
\]
$\theta = 30^\circ$.

\textbf{Final Answer:} $30^\circ$.

\hrulefill

\subsection*{Problem 20: SCalcET9 12.4.AE.001}
\textbf{Given:} $\vec{a} = \langle 2, 2, 2 \rangle, \vec{b} = \langle 2, 5, -4 \rangle$.
Fill in the blanks.
\[
\vec{a} \times \vec{b} = \begin{vmatrix} \ihat & \jhat & \khat \\ 2 & 2 & 2 \\ 2 & 5 & -4 \end{vmatrix}
\]
Row 1 expansion:
\[
= \begin{vmatrix} 2 & 2 \\ 5 & -4 \end{vmatrix}\ihat - \begin{vmatrix} 2 & 2 \\ 2 & -4 \end{vmatrix}\jhat + \begin{vmatrix} 2 & 2 \\ 2 & 5 \end{vmatrix}\khat
\]
Calculate:
$i: (-8 - 10) = -18$.
$j: (-8 - 4) = -12$. With outer minus: $+12$.
$k: (10 - 4) = 6$.

Filling the blanks in the PDF structure:
First line blanks: $2, 2$ (in j det), $2, 2$ (in k det).
Second line: $(-8 - 10)\ihat - (-8 - 4)\jhat + (10 - 4)\khat$.
Result: $-18\ihat + 12\jhat + 6\khat$.

\textbf{Final Answer:} $\langle -18, 12, 6 \rangle$.

\newpage

\section{Part 3: In-Depth Analysis of Problems and Techniques}

\subsection{A) Problem Types and General Approach}
\begin{itemize}
    \item \textbf{Type 1: Direct Calculation (Problems 1-5, 20):} These require setting up the $3 \times 3$ determinant. \textit{Strategy:} Always write $\ihat, \jhat, \khat$ in the top row. Be extremely careful with the middle term's negative sign.
    \item \textbf{Type 2: Algebraic Properties (Problems 6, 7):} Using the distributivity of the cross product. \textit{Strategy:} Treat it like FOILing polynomials, but respect the order ($\vec{u} \times \vec{v} = -\vec{v} \times \vec{u}$) and remember $\vec{u} \times \vec{u} = \vec{0}$.
    \item \textbf{Type 3: Conceptual/Meaningfulness (Problem 8):} Testing definitions. \textit{Strategy:} Check inputs/outputs. Cross product needs two vectors; dot product needs two vectors.
    \item \textbf{Type 4: Geometric Applications (Problems 9, 10, 19):} Using magnitude ($|\vec{a}||\vec{b}|\sin \theta$) and direction (RHR). \textit{Strategy:} Use $\tan \theta$ relationship when dot and cross products are both known.
    \item \textbf{Type 5: Area and Volume (Problems 13-18):} \textit{Strategy:} Parallelogram Area = $|\vec{a} \times \vec{b}|$. Triangle Area = $\frac{1}{2}|\vec{a} \times \vec{b}|$. Parallelepiped Volume = Scalar Triple Product. Coplanar check = Volume is 0.
\end{itemize}

\subsection{B) Key Algebraic and Calculus Manipulations}
\begin{enumerate}
    \item \textbf{The Pythagorean Trig Identity (Problem 4):} The simplification required $\sin^2 t + \cos^2 t = 1$ to verify orthogonality. Without this, the dot product check would not explicitly equal zero.
    \item \textbf{Anticommutativity (Problem 11):} Recognizing $\vec{b} \times \vec{a} = -(\vec{a} \times \vec{b})$ saved time by avoiding a second full determinant calculation.
    \item \textbf{Ratio of Magnitude to Dot Product (Problem 19):} Using $\frac{|\vec{a} \times \vec{b}|}{\vec{a} \cdot \vec{b}} = \frac{AB \sin \theta}{AB \cos \theta} = \tan \theta$. This is a powerful trick to find angles without finding vectors $\vec{a}$ or $\vec{b}$ explicitly.
\end{enumerate}

\section{Part 4: "Cheatsheet" and Tips for Success}

\subsection{Summary of Formulas}
\begin{itemize}
    \item $\vec{a} \times \vec{b} = \langle a_2b_3 - a_3b_2, -(a_1b_3 - a_3b_1), a_1b_2 - a_2b_1 \rangle$
    \item $|\vec{a} \times \vec{b}| = |\vec{a}| |\vec{b}| \sin \theta$ (Area of Parallelogram)
    \item Volume = $|\vec{a} \cdot (\vec{b} \times \vec{c})|$
    \item $\vec{a} \cdot (\vec{a} \times \vec{b}) = 0$ (Always orthogonal)
\end{itemize}

\subsection{Tricks and Pitfalls}
\begin{itemize}
    \item \textbf{The "Middle Minus":} The most common student error is forgetting the negative sign on the $\jhat$ component in the determinant expansion.
    \item \textbf{Order Matters:} $\vec{a} \times \vec{b} \neq \vec{b} \times \vec{a}$. If you swap the order, you MUST flip the sign.
    \item \textbf{Parallel Vectors:} If $\vec{a} \times \vec{b} = \vec{0}$, the vectors are parallel ($\theta = 0$ or $180^\circ$).
    \item \textbf{2D Cross Product:} To find the area of a polygon in 2D (z=0), just fill the z-column with 0s. The result will always be $\langle 0, 0, k \rangle$. The area is $|k|$.
\end{itemize}

\newpage

\section{Part 5: Conceptual Synthesis and The "Big Picture"}

\subsection{A) Thematic Connections}
The central theme of this topic is \textbf{Orthogonality in Space}. Just as the derivative allows us to find tangent lines (linear approximations), the cross product allows us to construct \textbf{normal vectors}. A normal vector is the key to defining a plane. This connects to linear algebra (determinants) and geometry (area/volume).

\subsection{B) Forward and Backward Links}
\begin{itemize}
    \item \textbf{Backward:} It relies on the Dot Product (orthogonality check) and Determinants.
    \item \textbf{Forward:}
    \begin{itemize}
        \item \textbf{Equations of Planes:} To write the equation of a plane, you need a normal vector. You get this by crossing two vectors that lie in the plane.
        \item \textbf{Surface Integrals (Calc 3):} Later, when calculating flux (flow across a surface), you will calculate $\vec{r}_u \times \vec{r}_v$ to find the normal vector to a curved surface. This is the bedrock of Stokes' Theorem and the Divergence Theorem.
    \end{itemize}
\end{itemize}

\section{Part 6: Real-World Application and Modeling}

\subsection{A) Concrete Scenarios}
\begin{enumerate}
    \item \textbf{Financial Volatility Surface Modeling:} In quantitative finance, options pricing relies on "volatility surfaces" (implied volatility plotted against strike price and time to maturity). To analyze the "skew" or local steepness of this surface for hedging purposes, quants treat the surface as a geometric object. The normal vector (found via cross product of tangent vectors along the strike and time axes) indicates the direction of maximum sensitivity.
    \item \textbf{Physics (Torque):} If you apply a force $\vec{F}$ to a wrench at a position $\vec{r}$ from the pivot, the torque (turning force) is $\vec{\tau} = \vec{r} \times \vec{F}$.
\end{enumerate}

\subsection{B) Model Problem Setup: Volatility Surface Normal}
\textbf{Scenario:} A quantitative analyst models an implied volatility surface $\sigma(k, t)$ where $k$ is strike price and $t$ is time.
\textbf{Variables:}
\begin{itemize}
    \item Position vector $\vec{P}(k, t) = \langle k, t, \sigma(k, t) \rangle$.
    \item Tangent along strike: $\vec{T}_k = \frac{\partial \vec{P}}{\partial k} = \langle 1, 0, \frac{\partial \sigma}{\partial k} \rangle$.
    \item Tangent along time: $\vec{T}_t = \frac{\partial \vec{P}}{\partial t} = \langle 0, 1, \frac{\partial \sigma}{\partial t} \rangle$.
\end{itemize}
\textbf{The Model:} To find the "Flux" of market risk or to define the plane tangent to the current market state, the analyst must solve for the normal vector $\vec{N}$:
\[
\vec{N} = \vec{T}_k \times \vec{T}_t = \begin{vmatrix} \ihat & \jhat & \khat \\ 1 & 0 & \sigma_k \\ 0 & 1 & \sigma_t \end{vmatrix} = \langle -\sigma_k, -\sigma_t, 1 \rangle
\]

\newpage

\section{Part 7: Common Variations and Untested Concepts}

\subsection{Missing Concept 1: The Triple Vector Product}
The homework covered the \textit{Scalar} Triple Product ($\vec{a} \cdot (\vec{b} \times \vec{c})$), but missed the \textit{Vector} Triple Product: $\vec{a} \times (\vec{b} \times \vec{c})$.
\textbf{Formula (BAC-CAB Rule):} $\vec{a} \times (\vec{b} \times \vec{c}) = \vec{b}(\vec{a} \cdot \vec{c}) - \vec{c}(\vec{a} \cdot \vec{b})$.
\textbf{Example:} If $\vec{a}=\ihat, \vec{b}=\jhat, \vec{c}=\ihat$, then $\ihat \times (\jhat \times \ihat) = \ihat \times (-\khat) = \jhat$.
Using formula: $\jhat(\ihat \cdot \ihat) - \ihat(\ihat \cdot \jhat) = \jhat(1) - \ihat(0) = \jhat$. Matches.

\subsection{Missing Concept 2: Distance from Point to Line}
The distance $d$ from a point $P$ to a line passing through $Q$ with direction $\vec{u}$ is given by $d = \frac{|\vec{QP} \times \vec{u}|}{|\vec{u}|}$.
\textbf{Example:} Distance from $P(1,1,1)$ to line through origin with direction $\langle 1,0,0 \rangle$.
$\vec{QP} = \langle 1,1,1 \rangle$, $\vec{u}=\langle 1,0,0 \rangle$.
$\vec{QP} \times \vec{u} = \langle 0, 1, -1 \rangle$. Magnitude $\sqrt{2}$.
$|\vec{u}| = 1$. Distance = $\sqrt{2}$.

\section{Part 8: Advanced Diagnostic Testing: "Find the Flaw"}

\subsection{Problem 1: Commutativity Error}
\textbf{Statement:} Find $\vec{a} \times \vec{b}$ given $\vec{a}=\langle 1,0,0 \rangle, \vec{b}=\langle 0,1,0 \rangle$.
\textbf{Flawed Solution:} Since the cross product is the area, and area is independent of order, $\vec{a} \times \vec{b} = \vec{b} \times \vec{a}$. We calculate $\vec{b} \times \vec{a} = -\khat$. Therefore $\vec{a} \times \vec{b} = -\khat$.
\textbf{The Error:} The cross product is anti-commutative, not commutative.
\textbf{Correction:} $\vec{a} \times \vec{b} = -(\vec{b} \times \vec{a}) = -(-\khat) = \khat$.

\subsection{Problem 2: Determinant Sign Error}
\textbf{Statement:} Calculate $\langle 1, 2, 3 \rangle \times \langle 4, 5, 6 \rangle$.
\textbf{Flawed Solution:}
$\ihat(12-15) + \jhat(6-12) + \khat(5-8) = -3\ihat - 6\jhat - 3\khat$.
\textbf{The Error:} The student forgot the negative sign attached to the $\jhat$ component expansion.
\textbf{Correction:} $\ihat(-3) - \jhat(-6) + \khat(-3) = -3\ihat + 6\jhat - 3\khat$.

\subsection{Problem 3: Distributing Dot over Cross}
\textbf{Statement:} Simplify $\vec{a} \cdot (\vec{a} \times \vec{b})$.
\textbf{Flawed Solution:} Distribute the dot product: $(\vec{a} \cdot \vec{a}) \times (\vec{a} \cdot \vec{b}) = |\vec{a}|^2 \times (\vec{a} \cdot \vec{b})$. This is meaningless.
\textbf{The Error:} The dot product does not distribute over the cross product; this is a geometric identity problem.
\textbf{Correction:} The vector $\vec{a} \times \vec{b}$ is orthogonal to $\vec{a}$. Therefore, the dot product is exactly 0.

\subsection{Problem 4: Scalar vs Vector in Triple Product}
\textbf{Statement:} Evaluate $(\ihat \times \jhat) \times \khat$.
\textbf{Flawed Solution:} $\ihat \times \jhat = 1$ (Area is 1). Then $1 \times \khat = \khat$.
\textbf{The Error:} $\ihat \times \jhat$ results in the vector $\khat$, not the scalar number 1.
\textbf{Correction:} $(\ihat \times \jhat) \times \khat = \khat \times \khat = \vec{0}$.

\subsection{Problem 5: Incorrect Triangle Area Formula}
\textbf{Statement:} Find area of triangle with vectors $\vec{a}, \vec{b}$.
\textbf{Flawed Solution:} Area $= |\vec{a} \times \vec{b}|$.
\textbf{The Error:} $|\vec{a} \times \vec{b}|$ is the area of the parallelogram; the triangle is half of that.
\textbf{Correction:} Area $= \frac{1}{2}|\vec{a} \times \vec{b}|$.

\end{document}